\documentclass[10pt,a4paper]{amsart}
\usepackage[margin=19mm]{geometry}
\usepackage{amsmath,amssymb,mathtools}
\usepackage[scr=rsfs,cal=euler]{mathalfa}
\usepackage{xfrac}
\usepackage[unicode,hidelinks,bookmarks=false]{hyperref}
\newcommand{\ie}{i.e.\ }
\newcommand{\cf}{cf.\ }
\newcommand{\resp}{respectively\ }
\newcommand{\Subsection}{\subsection}
\providecommand{\nobreakdash}{\nobreak\hbox{-}}
\setlength{\emergencystretch}{2em}
\multlinegap0pt
\usepackage{xcolor}
\usepackage{amssymb}
\newtheorem{proposition}{Proposition}
\newtheorem{lemma}{Lemma}
\newtheorem{corollary}{Corollary}
\newcommand{\R}{\mathbb{R}}
\newcommand{\Rmin}{\widetilde{\mathcal{R}}}
\newcommand{\SO}{\mathrm{SO}}
\newcommand{\Stab}{\mathrm{Stab}}
\newcommand{\Uc}{\mathcal{U}}
\newcommand{\sgn}{\mathrm{sign}}
\title{Linguistic Holonomy and Statistical Watermarks:\\ Inner Geometry of Meaning-Preserving Transformations}
\author{Daniele Corradetti}
\date{Source: arXiv:2608.19369v1 (CC BY 4.0)}
\begin{document}
\maketitle

\noindent\emph{Source and licence.} Linguistic Holonomy and Statistical Watermarks:\\ Inner Geometry of Meaning-Preserving Transformations. Version arXiv:2608.19369v1 (CC BY 4.0), distributed by arXiv under the Creative Commons Attribution 4.0 International licence, http://creativecommons.org/licenses/by/4.0/, as recorded in the arXiv licence metadata for this exact version and shown on the abstract page https://arxiv.org/abs/2608.19369v1. Full licence notice: \texttt{licenses/arxiv-2608.19369-CC-BY-4.0.txt}.\par\smallskip

\noindent\textit{Editorial scope.} Counted unit: the article's proposition prop:holonomy (source0.tex lines 481--503). The article's complete original proof of it is delivered with it (source0.tex lines 505--539, one \texttt{proof} environment of 35 lines). No mathematics is written, repaired or re-derived: the statement and the proof are the article's own lines, in the article's own notation, and no retained line is substituted or re-flowed.  Retained uncounted as the source-local context the proof needs: lines 265--267; lines 288--292; lines 295--297; lines 345--360; lines 453--459; lines 549--553.  Nothing was omitted from any retained span, so the package carries no omission record.  No retained line contains a citation, so the wrapper carries no bibliography and no bibliography entry.  No retained line contains a figure environment, an image-inclusion command or drawing code, so no image is delivered and the package declares no asset dependency.  Editorial formatting only, in addition: an AMS \texttt{amsart} preamble that restates the article's own declarations and macros used by the retained lines, namely the environments \texttt{proposition}, \texttt{lemma}, \texttt{corollary} on separate counters, together with the standard packages those lines need.  The retained environments print in this document as Proposition 1, Lemma 1, Corollary 1 in order of appearance, whereas the article prints the same items with the numbering of its own full document.  The complete native source of the article as distributed in the arXiv e-print for this exact version is bundled unchanged as \texttt{sources/208083/source0.tex}.
\begin{equation}\label{eq:cosine}
d_{*,\cos}(x,y) = 1 - \frac{x\cdot y}{\|x\|\,\|y\|}.
\end{equation}

\begin{equation}\label{eq:chain}
\psi(\Uc(\lambda)) = \{v_0, v_1, \dots, v_L\}, \qquad
v_0 = \psi(\lambda), \quad
v_i = \psi\bigl((U_i \circ \cdots \circ U_1)(\lambda)\bigr).
\end{equation}

\begin{equation}\label{eq:deficit}
\delta_\Uc(\lambda) := d_*(v_0, v_L),
\end{equation}

\begin{proposition}[Degeneracy of the signature]\label{prop:sig}
Let $R \in \SO(n)$ with non-trivial rotation angles
$\theta_1,\dots,\theta_p \in (0,\pi]$, and set
$M = I_n - \tfrac12(R+R^T)$. Then $M$ is positive semidefinite, its
spectrum is
\begin{equation}\label{eq:spectrum}
\bigl\{\, 1 - \cos\theta_k \ \text{with multiplicity } 2 \,\bigr\}_{k=1}^{p}
\ \cup\ \bigl\{\, 0 \ \text{with multiplicity } n-2p \,\bigr\},
\end{equation}
and consequently
\begin{equation}\label{eq:signature}
\sgn(M) = (2p,\, 0,\, n-2p).
\end{equation}
The negative index vanishes identically, $\mathrm{rank}(M) = 2p$, and
the signature is a function of the single integer $p$.
\end{proposition}

\begin{lemma}[The minimal rotation is parallel transport]\label{lem:transport}
Let $\hat{x},\hat{y}$ be distinct, non-antipodal unit vectors of $\R^n$
and $P = \mathrm{span}(\hat{x},\hat{y})$. Then $\Rmin^{x,y}$ is the
parallel transport of the round metric of $S^{n-1}$ along the minimising
geodesic from $\hat{x}$ to $\hat{y}$, extended to $\R^n$ as the ambient
rotation fixing $P^{\perp}$ pointwise.
\end{lemma}

\begin{corollary}[Gauss--Bonnet]\label{cor:gaussbonnet}
For $n = 3$ and $L = 2$ the holonomy is the rotation of the tangent
plane at $\hat{v}_0$ by the spherical excess of the geodesic triangle
$\hat{v}_0\hat{v}_1\hat{v}_2$, that is by its area.
\end{corollary}

\begin{proposition}[Endpoint--path factorisation]\label{prop:holonomy}
Let $V = (v_0,\dots,v_L)$ be a chain as in \eqref{eq:chain}, with
$n \ge 3$ and no two consecutive states antipodal. Put
$\mathrm{R}_{\mathrm{dir}} = \Rmin^{v_0,v_L}$ and
\begin{equation}\label{eq:holdef}
H := \mathrm{R}_{\mathrm{dir}}^{-1}\,\mathrm{R}_\Uc .
\end{equation}
Then
\begin{enumerate}
\item[(i)] $H\hat{v}_0 = \hat{v}_0$, so that
$H \in \Stab(\hat{v}_0) \cong \SO(n-1)$, and
$\mathrm{R}_\Uc = \mathrm{R}_{\mathrm{dir}}H$ canonically;
\item[(ii)] $\delta_\Uc = 1 - \langle \hat{v}_0,
\mathrm{R}_{\mathrm{dir}}\hat{v}_0\rangle$ depends only on
$\mathrm{R}_{\mathrm{dir}}$, and is therefore blind to $H$;
\item[(iii)] $H$ is the Riemannian holonomy, based at $\hat{v}_0$, of
the closed geodesic polygon obtained by closing the path with the
geodesic from $\hat{v}_L$ back to $\hat{v}_0$;
\item[(iv)] for every $u \in S^{n-1}$ and every
$H_0 \in \Stab(\hat{v}_0)$ there is a chain with initial state
$\hat{v}_0$, final state $u$ and holonomy exactly $H_0$.
\end{enumerate}
\end{proposition}

\begin{proof}
(i) By construction $\Rmin^{v_{i-1},v_i}\hat{v}_{i-1} = \hat{v}_i$, so
by induction $\mathrm{R}_\Uc \hat{v}_0 = \hat{v}_L$; and
$\mathrm{R}_{\mathrm{dir}}\hat{v}_0 = \hat{v}_L$ by definition. Since
$\mathrm{R}_{\mathrm{dir}}$ is orthogonal,
$H\hat{v}_0 = \mathrm{R}_{\mathrm{dir}}^{T}\hat{v}_L = \hat{v}_0$. An
element of $\SO(n)$ fixing a unit vector preserves its orthogonal
hyperplane and restricts there to an element of $\SO(n-1)$.

(ii) Immediate from \eqref{eq:deficit}, \eqref{eq:cosine} and
$\hat{v}_L = \mathrm{R}_{\mathrm{dir}}\hat{v}_0$.

(iii) By Lemma~\ref{lem:transport}, $\mathrm{R}_\Uc$ and
$\mathrm{R}_{\mathrm{dir}}$ are the transports along the two paths, and
\eqref{eq:holdef} is the transport around the closed circuit.

(iv) It suffices to realise every $H_0$ as the holonomy of a closed
geodesic polygon based at $\hat{v}_0$: appending the single geodesic leg
from $\hat{v}_0$ to $u$ multiplies the loop matrix on the left by
$\Rmin^{\hat{v}_0,u}$ and leaves the holonomy unchanged. Fix a two-plane
$Q \subset \hat{v}_0^{\perp}$ and an angle $\alpha \in (0,2\pi)$, and let
$\Sigma = S^{n-1}\cap(\mathrm{span}(\hat{v}_0)\oplus Q)$, a totally
geodesic two-sphere through $\hat{v}_0$ whose tangent space there is
$Q$. A geodesic triangle in $\Sigma$ with vertex $\hat{v}_0$ and area
$\alpha$ exists, since the area of a geodesic triangle on the unit
two-sphere sweeps the whole of $(0,2\pi)$; transport around a loop
contained in a totally geodesic submanifold is the transport computed
inside it, extended by the identity on the normal directions, so by
Corollary~\ref{cor:gaussbonnet} below the holonomy is the rotation of
$Q$ by $\alpha$ and the identity elsewhere. By the normal form used in
Proposition~\ref{prop:sig}, every element of $\SO(n-1)$ is a product of
at most $\lfloor (n-1)/2 \rfloor$ such plane rotations, and the
corresponding polygons, each beginning and ending at $\hat{v}_0$, may be
concatenated.
\end{proof}

\end{document}
